\documentclass[11pt]{article}
\usepackage[T1]{fontenc}
\usepackage[utf8]{inputenc}
\usepackage{lmodern,microtype}
\usepackage[a4paper,margin=25mm]{geometry}
\usepackage{amsmath,amssymb,amsthm,mathtools,booktabs,array}
\usepackage{xcolor,enumitem,listings}
\usepackage[hidelinks,breaklinks]{hyperref}
\usepackage{fancyhdr}
\definecolor{ink}{HTML}{17374C}
\definecolor{muted}{HTML}{526473}
\definecolor{pale}{HTML}{F3F6F8}
\hypersetup{pdftitle={One computation axiom for iPRA},pdfauthor={Research report prepared for Vladimir Reshetnikov}}
\pagestyle{fancy}\fancyhf{}
\fancyhead[L]{\small\color{muted}One computation axiom for iPRA}
\fancyhead[R]{\small\color{muted}11 November 2026}
\fancyfoot[C]{\thepage}\setlength{\headheight}{14pt}
\setlength{\parindent}{0pt}\setlength{\parskip}{5pt}
\setlist{itemsep=3pt,topsep=5pt}
\lstset{basicstyle=\ttfamily\small,breaklines=true,columns=fullflexible,keepspaces=true,showstringspaces=false,backgroundcolor=\color{pale},frame=single,rulecolor=\color{pale}}
\newtheorem{theorem}{Theorem}[section]
\newtheorem{lemma}[theorem]{Lemma}
\newtheorem{corollary}[theorem]{Corollary}
\newcommand{\IPRA}{\mathsf{iPRA}}
\newcommand{\N}{\mathcal N}
\newcommand{\Zsix}{\mathcal Z_6}
\newcommand{\Zseven}{\mathcal Z_7}
\newcommand{\Q}{\mathsf Q}\newcommand{\D}{\mathsf D}
\newcommand{\nf}{\operatorname{nf}}\newcommand{\It}{\widehat I}
\newcommand{\Eqv}{\dashv\vdash}
\newcommand{\code}[1]{\nolinkurl{#1}}
\begin{document}
\begin{titlepage}
\vspace*{10mm}
{\large\color{muted}CONCISE INTUITIONISTIC ARITHMETIC\par}
\vspace{8mm}
{\Huge\bfseries\color{ink}One computation axiom\\[4pt]for iPRA\par}
\vspace{8mm}
{\LARGE A six-rule presentation with zero-normal terms\par}
\vspace{10mm}
{\large A checked simplification of the current ConciseIPRA proposal\par}
\vspace{13mm}
\begin{minipage}{0.94\textwidth}
\textbf{Result.}
Make only \(I(0,a;u.h)=a\) definitional. Use a term layout in which the
initializer \(a\) stays outside the state binder. The existing
single-introduction calculus then needs six primitive rules, including
only one computation schema. The complete, strictly validating tagged
Python checker decreases from 3,145 to 3,095 normalized bytes, from
1,024 to 1,013 significant tokens, and from 1,080 to 1,075 AST nodes.
The normalization and its input checks are included.

\medskip
\textbf{Equivalence.}
The new theory is bi-interpretable and equiconsistent with full
first-order intuitionistic primitive recursive arithmetic. The numerical
domains and equality are unchanged. Both interpretation composites are
internally the identity; on old term syntax this uses \(t=\nf(t)\).

\medskip
\textbf{Scope.}
This is a local reduction with a written proof, a working checker,
proof translations, and compiled Lean verification of the normalization
bridge. The entire six-rule equivalence is not formalized in Lean.
The construction does not beat the repository's 521-byte arity-coded
checker and does not establish an absolute minimum. Earlier packages
explicitly proposed the underlying idea; the present contribution
makes it work at lower cost in the direct tagged layout.
\end{minipage}
\vfill
{\large Prepared for Vladimir Reshetnikov\par}
{11 November 2026\par}\vspace{4mm}
{\small Baseline: ProveIt commit\\
\code{ecb3b847290c33b1ef7f04571cac63e97c73a54e}\par}
\end{titlepage}
\begin{abstract}
The native ConciseIPRA calculus uses numerical iterator terms, equations,
and two binders \(\Q,\D\). Its current reductions already remove separate
logical axioms and combine the case and induction introductions. We
complete an earlier suggested reduction: use the normal forms of the
single erasing rewrite \(I(0,a;u.h)\to a\) as terms. Hereditary substitution
then absorbs zero computation, leaving successor computation as the sole
nonlogical axiom schema. A different closed representative of zero,
\(K=I(1,0;u.0)\), preserves the single-introduction construction.
We give all six rules, effective translations in both directions, both
round trips, and the transfer of the existing full iPRA bi-interpretation.
The changed syntax, scope validation, and hereditary substitution are
charged in a matched implementation comparison.
\end{abstract}
\setcounter{tocdepth}{1}
\tableofcontents\newpage

\section{What changes, and what was already known}
The comparison is to the current repository at the commit on the title
page, including its fourth delivery of research packages. The native
kernel has eight constructors: hypothesis, two computation equations,
equational transport, \(\Q\) elimination, \(\D\) elimination,
\(\D\) introduction by cases, and equational \(\Q\) introduction by
progress. The existing \code{witness_rule}/\code{NativeOneIntro}
development combines the last two introductions. The existing
\code{one_context} proposal replaces two-template equality transport
by one crossed numerical context \cite{readme,native,oneintro,onecontext}.

The present change is
\[
 \boxed{\text{Remove zero computation from proofs; normalize it in term syntax.}}
\]
It helps only if the normalization costs less than the deleted
proof-checking case. Section~\ref{sec:size} measures the complete cost,
including the canonical-term check.

\subsection{Direct antecedents}
Definitional zero computation is not a new mathematical idea.
The subsection \emph{Making zero computation definitional} in
\code{prevalidated_elimination/theory/research.tex} already proposes it
and explains why the examined implementation lost: its arity codec
places an initializer under a dummy slot, so erasing an iterator
requires lowering coordinates. \code{uniform_zero/report.tex}
explicitly asks whether a different binder layout could make the
conversion cheaper \cite{priorzero,uniformzero}.

The direct tagged syntax has that different layout. In \(I(c,a;u.h)\),
only \(h\) binds \(u\); \(c,a\) remain at the surrounding scope.
Erasing a zero-counter iterator returns \(a\) without scope adjustment.
This report supplies the smaller implementation, the revised witness
construction, and the proof obligations for both interpretation
composites.

\subsection{The exact simplicity claim}
The result has two independently checkable parts:
\begin{enumerate}
\item The complete axiomatization has six primitive proof rules and only
one computation axiom schema.
\item Against an equally factored seven-rule tagged control, the standalone
checker is smaller in normalized bytes, significant tokens, and AST nodes.
Both programs validate numerical scopes and formula sorts. The public
term grammars differ intentionally.
\end{enumerate}
The 521-byte programs use a different arity-coded protocol, guarded
public language, and auxiliary resources. They remain smaller programs.
No global lower bound, absolute optimum, or new record over that
frontier is claimed.

\section{Term syntax and hereditary substitution}\label{sec:syntax}
\subsection{Raw terms and normal forms}
Start with raw terms
\[
 t ::= x\mid0\mid S(t)\mid I(c,a;u.h),
\]
where only \(u\) is bound, and only in \(h\). There is one numerical sort.
Define the bottom-up operation
\begin{align*}
 \nf(x)&=x,&\nf(0)&=0,&\nf(St)&=S\nf(t),\\
 \nf(I(c,a;u.h))&=
 \begin{cases}
 \nf(a),&\nf(c)\text{ is literally }0,\\
 I(\nf(c),\nf(a);u.\nf(h)),&\text{otherwise}.
 \end{cases}
\end{align*}
New terms are exactly the fixed points of \(\nf\): well-scoped terms
with no literal-zero-counter iterator at any depth.
For already normal arguments define the smart constructor
\[
 \It(c,a;u.h)=
 \begin{cases}a,&c\equiv0,\\ I(c,a;u.h),&c\not\equiv0.\end{cases}
                                                        \tag{SC}\label{eq:smart}
\]
Here \(\equiv\) is literal syntax. This constructs syntax; it adds
no object-language arithmetic function or axiom.

Only zero computation is normalized. In particular
\[
                   K:=I(S0,0;u.0)                       \tag{K}\label{eq:alias}
\]
is a normal closed term distinct in syntax from \(0\). Its equality
to zero is proved using the remaining axiom.

\subsection{Substitution and formulas}
All substitutions are capture-avoiding substitutions followed by
normalization:
\[
                    t[\sigma]_0:=\nf(t[\sigma]).
                                                        \tag{HS}\label{eq:hsub}
\]
Compute this directly by recursively rebuilding with \(\It\).
Substitute in the iterator counter and initializer at the current scope,
and use the lifted substitution only in its body.
Renaming preserves normality and is ordinary renaming.

Formulas are
\[
                 A::=(s=t)\mid\Q_y A\mid\D_{x,y}A .
\]
The two binders bind the displayed numerical variables. In de Bruijn
notation \(y\) has index \(0\) and \(x\) index \(1\) in a \(\D\)-matrix.
Formula substitution normalizes its numerical terms and lifts under
each binder. Normalization never changes a formula's root constructor.

Every \(B(t,r)\) below means simultaneous hereditary instantiation
of \(B(x,y)\), and \(h(t)\) means \(h[t/u]_0\). Bound variables are
renamed fresh before substitution. Neither the normalizer nor the
introduction rule tests whether an arbitrary term \emph{denotes} zero.

\section{The complete six-rule theory}\label{sec:rules}
Let \(\Zsix\) have finite lists \(\Gamma\) of reusable assumptions.
Its designated contradiction is \(0=S0\). These are all six primitive
rules; IS is a zero-premise axiom schema.

\paragraph{1. Hypothesis.}
\[
             \frac{\ }{\Gamma\vdash A}\quad(A\in\Gamma).
                                                        \tag{Hyp}
\]
\paragraph{2. Successor computation.}
For arbitrary normal \(c,a,h\),
\[
 \frac{\ }{\Gamma\vdash
 I(Sc,a;u.h)=h[\It(c,a;u.h)/u]_0}.
                                                        \tag{IS}\label{eq:is}
\]
The left side is normal because \(Sc\) is not literal zero.

\paragraph{3. One-context crossed equality.}
For a normal numerical template \(h(v)\),
\[
 \frac{\Gamma\vdash s=t\qquad\Gamma\vdash h(s)=r}
      {\Gamma\vdash r=h(t)}.
                                                        \tag{X}\label{eq:x}
\]
There is no restriction on \(s\).

\paragraph{4. Closure elimination.}
\[
 \frac{\Gamma\vdash\Q_y B(y)\qquad\Gamma\vdash B(0)}
      {\Gamma\vdash B(t)}.
                                                        \tag{QE}
\]
\paragraph{5. Witness elimination.}
\[
 \frac{\Gamma\vdash\D_{x,y}B(x,y)
       \qquad\Gamma,\Q_y B(b,y)\vdash A}
      {\Gamma\vdash A}.
                                                        \tag{DE}
\]
The eigenvariable \(b\) is fresh for \(\Gamma,A\) and the fixed matrix
parameters. It is not a freely chosen witness term.

\paragraph{6. Single introduction.}
\[
 \frac{\Gamma,B(t,\eta)\vdash B(t,Sz)}
      {\Gamma\vdash\D_{x,y}B(x,y)},\qquad
 \eta=\begin{cases}
 z,&B\text{ is a literal equation and }t\equiv0,\\
 0,&\text{otherwise}.
 \end{cases}                                            \tag{DI}\label{eq:di}
\]
The variable \(z\) is fresh for \(\Gamma,t\) and the fixed parameters
of \(B\). The displayed local assumption is discharged.

There are no additional primitive reflexivity, symmetry, substitution,
cut, induction, successor, coverage, separation, or logical axioms.
The ordinary readings established by the interpretation are
\[
 \Q_yB(y)\rightsquigarrow B(0)\to\forall y\,B(y),\qquad
 \D_{x,y}B\rightsquigarrow
    \exists x\bigl(B(x,0)\to\forall y\,B(x,y)\bigr).
                                                        \tag{Reading}
\]
These readings explain the binders and do not add proof rules.

\section{Recovering equality and both introductions}\label{sec:derived}
\subsection{Reflexivity without zero computation}
\begin{lemma}\label{lem:rfl}
For every normal term \(r\), \(\Zsix\) proves \(r=r\).
\end{lemma}
\begin{proof}
IS with \(c=0,a=0\) and the constant body \(r\), weakened under the
state binder, gives \(L_r=r\), where \(L_r=I(S0,0;u.r)\).
Use this equation as both premises of X with identity template.
The result is \(r=r\).
\end{proof}

Symmetry follows from X with major \(s=s\), minor \(s=t\), and identity
template. Given \(r=s\) and \(s=t\), reverse the first equation and apply X
with major \(s=t\), minor \(s=r\), and the identity template; its
conclusion is \(r=t\). Congruence follows from X with major \(s=t\)
and reflexive minor \(h(s)=h(s)\).
Two X steps recover the full two-template transport:
\[
 \begin{gathered}
 s=t,\ a(s)=b(s)\ \Longrightarrow\ b(s)=a(t),\\
 s=t,\ b(s)=a(t)\ \Longrightarrow\ a(t)=b(t).
 \end{gathered}                                         \tag{Tr}
\]

\subsection{Structural facts}
Weakening, exchange, contraction, variable renaming, and ordinary
assumption cut are admissible by direct induction on a proof.
Weakening and renaming preserve the formula root and the literal-zero
status of a DI witness. For cut, replace uses of the discharged
assumption by its proof, renaming eigenvariables apart. This does not
substitute a numerical term into the witness.

Arbitrary numerical substitution of six-rule proofs is established
later through a two-introduction calculus. It cannot be justified by
an unqualified constructor induction: substitution may turn a
nonliteral witness into literal zero.

\subsection{A seven-rule bridge}
Define \(\Zseven\) on the same normal syntax by replacing DI with
\[
 \frac{\Gamma,B(t,0)\vdash B(t,Sz)}
      {\Gamma\vdash\D_{x,y}B(x,y)}
 \quad(\mathrm{DI}_{\rm case})
\]
for arbitrary matrices and witnesses, and
\[
 \frac{\Gamma,E(z)\vdash E(Sz)}
      {\Gamma\vdash\Q_y E(y)}
 \quad(\mathrm{QI}_{\rm prog})
\]
for literal equations only. The corresponding freshness conditions
are retained. These seven rules are stable under numerical substitution.

\begin{theorem}\label{thm:sixseven}
\(\Zsix\) and \(\Zseven\) derive exactly the same sequents.
\end{theorem}
\begin{proof}
To recover case introduction in \(\Zsix\), only an equation \(E(x,y)\)
at witness \(0\) is missing. Given
\[
                  \Gamma,E(0,0)\vdash E(0,Sz),
\]
use \(K=I(S0,0;u.0)\). It is normal and syntactically nonzero, and
IS proves \(K=0\). Under \(E(K,0)\), equational transport gives
\(E(0,0)\). Cut in the supplied proof and transport its conclusion
to \(E(K,Sz)\), using \(0=K\). DI at \(K\) selects case mode and
concludes the desired \(\D_{x,y}E\).
This uses equation transport and proof-cut, not an assumed numerical
substitution theorem.

For progress introduction of \(E(y)\), add an unused witness \(x\)
and use DI at witness \(0\). It gives \(\D_{x,y}E(y)\).
DE exposes exactly the desired \(\Q_yE(y)\), independent of its
fresh witness; Hyp returns it.

Conversely, a case-mode DI is the case rule. For a progress-mode DI,
its matrix is an equation and its witness is zero.
\(\mathrm{QI}_{\rm prog}\) gives \(\Q_y B(0,y)\).
Under \(B(0,0)\), QE gives \(B(0,Sz)\); case introduction at zero
then produces \(\D_{x,y}B\).
\end{proof}

\begin{corollary}[Equation induction]\label{cor:ind}
For a literal equation \(E\), the rule
\[
 \frac{\Gamma\vdash E(0)\qquad\Gamma,E(z)\vdash E(Sz)}
      {\Gamma\vdash E(t)}
\]
is derived, with the usual freshness condition.
\end{corollary}
\begin{proof}
Use recovered progress introduction and QE.
\end{proof}
The new zero representative \(K\) matters: the old
\(I(0,0;u.0)\) has become the syntax \(0\) and cannot select case mode.

\section{Normalization algebra}\label{sec:nfalgebra}
\begin{lemma}[Normal forms]\label{lem:nf}
The operation \(\nf\) produces a normal term, is idempotent, and fixes
exactly normal terms. It preserves scopes and commutes with variable
renaming.
\end{lemma}
\begin{proof}
Use structural induction. Normalize all iterator children. If its
counter is zero, the initializer is already normal and has the
surrounding scope; otherwise the rebuilt iterator is normal.
Renaming cannot turn a variable, successor, or nonzero-headed iterator
into literal zero, so it commutes with the decision.
\end{proof}

The rewrite \(I(0,a;u.h)\to a\) terminates because it decreases the
number of nodes. A contraction in any position preserves \(\nf\):
it preserves a retained child's normal form and cannot affect the
result when it occurs in an erased body. Every term reduces to its
bottom-up normal form. Together these facts prove uniqueness of
normal forms, including contractions under binders.

\begin{lemma}[Substitution coherence]\label{lem:sub}
For every simultaneous raw substitution \(\sigma\),
\[
 \nf(t[\sigma])
       =\nf\bigl(\nf(t)[\,\nf\circ\sigma\,]\bigr).
                                                        \tag{NC}\label{eq:nfsub}
\]
The corresponding formula identity holds under both binders.
\end{lemma}
\begin{proof}
Induct on \(t\), with renaming compatibility for lifted substitutions.
If an iterator counter has \(\nf(c)=0\), the induction hypothesis
implies \(\nf(c[\sigma])=0\), so both sides return the same substituted
initializer. Otherwise substitution may or may not create zero, but
both sides use the same smart constructor on the same three
inductively equal children. The body uses the lifted substitution.
Formula induction uses one lift under \(\Q\) and two under \(\D\).
\end{proof}

\begin{corollary}
On normal terms, hereditary substitution obeys identity and composition and commutes
with weakening and binder instantiation.
\end{corollary}
\begin{proof}
Apply \eqref{eq:nfsub}, idempotence, and the raw substitution laws.
\end{proof}

The normalizer is a structural erasing pass. It does not run
successor iteration, test arithmetic equality, or evaluate arbitrary
primitive recursive programs. Substitution can enlarge a term by
copying its image, as in the baseline; the added normalization removes
nodes.

\section{The internal round trip}\label{sec:roundtrip}
Changing an iterator body is not automatically an instance of ordinary
one-hole term congruence. Pointwise equality of bodies must be lifted
through iteration.

\begin{lemma}[Iterator-body congruence]\label{lem:body}
In the raw native calculus, if \(h(u)=k(u)\) is a theorem with \(u\)
and passive parameters free, then
\[
                  I(n,a;u.h)=I(n,a;u.k)
\]
is a theorem.
\end{lemma}
\begin{proof}
Induct on \(n\) using literal-equation induction.
Both zero instances equal \(a\).
At a successor, the induction hypothesis equates the previous
values. Congruence through \(h\), followed by the pointwise theorem
instantiated at the second previous value, equates the updates.
Successor computation completes the step.
The pointwise equation is a theorem, not a retained hypothesis
allowed to vary with the iteration state.
\end{proof}
The existing \code{Q.iter_body_congr} and \code{Q.iter_congr} in
\code{RoundTripScalar.lean} formalize the ordinary scalar version;
the native/scalar bridge transfers it \cite{roundtrip,nativebridge}.

\begin{theorem}[Normalization is internally the identity]\label{thm:nfeq}
For every raw term \(t\), the raw native calculus proves
\[
                            t=\nf(t).                   \tag{NI}\label{eq:nfeq}
\]
For every raw formula \(A\), it proves \(A\Eqv\nf(A)\), uniformly
in its free numerical parameters.
\end{theorem}
\begin{proof}
Induct structurally on \(t\). Variables and zero use reflexivity;
successor uses congruence. In \(I(c,a;u.h)\), ordinary congruence
replaces \(c,a\) using their induction hypotheses, and
Lemma~\ref{lem:body} replaces \(h\).
If \(\nf(c)=0\), raw zero computation then yields \(\nf(a)\);
otherwise the term already is \(\nf(t)\).

For formulas, use these endpoint equalities at equations and
the derived native binder congruences at \(\Q,\D\).
Equivalently, use their ordinary intuitionistic readings and the
native/scalar equivalence. The matrix equivalences hold uniformly
in the bound variables, so implication and quantifier congruence
apply.
\end{proof}

This is an internal identity, not merely equality of standard-model
values. Inclusion of normal syntax into raw syntax preserves
substitution only up to this proved equality:
\[
 \N\vdash
 \operatorname{inc}(t[\sigma]_0)
 =\operatorname{inc}(t)[\operatorname{inc}\sigma].
\]
The left term description is definitionally the normal form of the
right; the displayed equality is a theorem. The two descriptions need
not be identical.

\section{Both proof translations}\label{sec:translations}
Let \(\N\) be the raw eight-constructor native kernel. Normalize
formulas by normalizing their terms and retaining their formula nodes.

\begin{theorem}\label{thm:oldnew}
If \(\Gamma\vdash_\N A\), then
\(\nf(\Gamma)\vdash_{\Zseven}\nf(A)\), hence also
\(\nf(\Gamma)\vdash_{\Zsix}\nf(A)\).
\end{theorem}
\begin{proof}
Translate the finite proof tree inductively.
Hypothesis membership is preserved.
A zero-computation leaf becomes reflexivity, supplied by
Lemma~\ref{lem:rfl}. A successor-computation leaf becomes IS by
Lemma~\ref{lem:sub}. Replace two-template transport by two X steps.
QE, DE, and both introductions commute with normalization by
\eqref{eq:nfsub} and its lifted forms. Progress matrices remain
equations. Normalization may erase free-variable occurrences but
never introduces them, so eigenvariable conditions are preserved.
Finally apply Theorem~\ref{thm:sixseven}.
\end{proof}

Do not normalize an old single-introduction proof and blindly retain
its introduction mode: a nonliteral witness can become literal zero.
The translation passes through the two-introduction systems.

\begin{theorem}\label{thm:newold}
If \(\Gamma\vdash_{\Zseven}A\), inclusion into raw syntax gives
\(\Gamma\vdash_\N A\). The same holds for \(\Zsix\).
\end{theorem}
\begin{proof}
First use Theorem~\ref{thm:sixseven} if necessary, then induct on the
seven-rule proof.
Hyp is unchanged. IS is raw successor computation followed by
replacing its endpoint with its proved normal form.
For X, use the native equalities between raw \(h(s),h(t)\) and their
hereditary instances, supplied by Theorem~\ref{thm:nfeq}, at both
ends of native one-context transport.

For QE, convert the normalized base to its raw instance, use native
QE, then convert the target to its normal instance.
For DE, the normal matrix itself is unchanged and renaming commutes
with normalization, so the fresh local closure has the right reading.
For case introduction, convert its raw local assumption \(B(t,0)\)
to its normal instance, cut in the translated premise, convert the
endpoint back to raw \(B(t,Sz)\), then use native case introduction.
For progress, use the corresponding coherence at the equation
instances and native equation progress.
All conversions are parameter-uniform; no eigenvariable enters the
retained context or external conclusion.
\end{proof}

\begin{corollary}[Faithfulness and composites]\label{cor:faithful}
For every raw sequent,
\[
 \Gamma\vdash_\N A
 \quad\Longleftrightarrow\quad
 \nf(\Gamma)\vdash_{\Zsix}\nf(A).
\]
Normal syntax followed by inclusion and normalization returns
literally itself. Raw syntax followed by normalization and inclusion
returns a formula internally equivalent to the original.
\end{corollary}
\begin{proof}
Use the two proof translations, Theorem~\ref{thm:nfeq}, and cut.
Idempotence proves the first composite and internal normalization
correctness proves the other.
\end{proof}

\section{Bi-interpretability with full iPRA}\label{sec:ipra}
Here \(\IPRA\) is full first-order intuitionistic primitive recursive
arithmetic: a symbol for each primitive recursive definition, its
defining equations, successor separation and injectivity, and
quantifier-free induction, in intuitionistic first-order logic with
equality. Only the induction schema is restricted to quantifier-free
formulas; the language of theorems is unrestricted.

The pinned repository proves the full native/iPRA bi-interpretation,
including both formula round trips, rather than only equality of
quantifier-free consequences \cite{nativebridge}.
Its intermediate theory \(T_{\rm eq}\) has intuitionistic first-order
logic over iterator terms, the two computation equations, numerical
coverage, \(0=S0\to\bot\), and equation induction. The inherited
arithmetic development derives predecessor, reflection of equality,
quantifier-free induction, pairing, and a compiler for primitive
recursion. The frozen Lean source chain is included.

Let \(J:\N\to\IPRA\) and \(H:\IPRA\to\N\) be those interpretations.
Define
\[
       J_0=J\circ\operatorname{inc},\qquad H_0=\nf\circ H .
\]
Both use the whole numerical domain and actual equality. No quotient
of numbers, restricted cut, or coded subclass replaces the domain;
only term descriptions are normalized.

\begin{theorem}[Main theorem]\label{thm:main}
\(\Zsix\) is bi-interpretable and equiconsistent with \(\IPRA\).
Both numerical composites have the identity graph \(x=y\).
\end{theorem}
\begin{proof}
Normalization and inclusion preserve and reflect consequence by
Section~\ref{sec:translations}, and their composites have the internal
identities of Corollary~\ref{cor:faithful}. Compose with the pinned
native/iPRA round trips. For every normal formula \(A\) and arbitrary
iPRA formula \(\varphi\), this yields
\[
 \Zsix\vdash A\leftrightarrow H_0J_0A,\qquad
 \IPRA\vdash\varphi\leftrightarrow J_0H_0\varphi ,
\]
where the first equivalence uses the derived native macros.
Variables are unchanged, so the numerical identity graph is \(x=y\).

Normalization fixes \(0,S0\). The native bridge identifies
\(0=S0\) with contradiction up to derivability. Thus
\[
                   \Zsix\vdash0=S0
           \quad\Longleftrightarrow\quad\IPRA\vdash\bot ,
\]
which proves equiconsistency.
\end{proof}

\subsection{Explicit logical macros}
All constructions in this subsection use smart iteration and
hereditary substitution. Put \(C(v,a,b)=\It(v,a;u.b)\), with \(u\)
fresh. Define an equation-to-formula selector \(L(v,s,t;B)\):
\begin{align*}
 L(v,s,t;r=q)&=(C(v,s,r)=C(v,t,q)),\\
 L(v,s,t;\Q_yB)&=\Q_yL(v,C(y,0,s),C(y,0,t);B),\\
 L(v,s,t;\D_{x,y}B)&=\D_{x,y}L(v,C(y,0,s),C(y,0,t);B).
\end{align*}
In every binder clause passive arguments are weakened into its scope
and bound names are chosen fresh.
Set \(G(v,A)=L(v,0,0;A)\), then define \(\operatorname{Sel}\)
recursively on its first formula:
\begin{align*}
 \operatorname{Sel}(v,s=t,B)&=L(v,s,t;B),\\
 \operatorname{Sel}(v,\Q_yA,B)&=
       \Q_y\operatorname{Sel}(v,A,G(y,B)),\\
 \operatorname{Sel}(v,\D_{x,y}A,B)&=
       \D_{x,y}\operatorname{Sel}(v,A,G(y,B)).
\end{align*}
These are finite recursive macros, not new primitive connectives.
The inherited \code{NativeGate.lean} proves the endpoint properties,
which transfer through normalization:
\[
 \operatorname{Sel}(0,A,B)\Eqv A,\qquad
 \operatorname{Sel}(Sz,A,B)\Eqv B,\qquad
 \vdash G(0,A),\quad G(Sz,A)\Eqv A .
\]
Every auxiliary bound variable below is fresh for the input formulas
and passive parameters, except for the variable intentionally bound by
the displayed existential or universal quantifier.
With \(p(y)=\It(y,0;u.C(u,1,0))\), define
\[
\begin{array}{rcl}
 \bot&:=&(0=S0),\\
 A\to B&:=&\Q_y\operatorname{Sel}(y,A,B),\\
 A\wedge B&:=&\Q_yG(y,\operatorname{Sel}(p(y),A,B)),\\
 \exists x\,A(x)&:=&\D_{x,y}G(y,A(x)),\\
 A\vee B&:=&\exists v\,\operatorname{Sel}(v,A,B),\\
 \forall y\,A(y)&:=&A(0)\wedge\Q_y A(y).
\end{array}
\]
The inherited \code{NativeLogic.lean} derives all intuitionistic
natural-deduction rules for these macros.
Excluded middle for arbitrary formulas is not added.
The result is full intuitionistic arithmetic, not only an equational
fragment.

\section{The standalone checker}\label{sec:checker}
\subsection{Strict external syntax}
The main checker is \code{code/zero_six.py}. It uses tagged JSON and
scoped de Bruijn indices:
\begin{center}
\begin{tabular}{ll}
\toprule
Syntax&JSON\\\midrule
variable \(i\)&\code{["v", i]}\\
\(0\)&\code{["z"]}\\
\(S(t)\)&\code{["s", t]}\\
\(I(c,a;u.h)\)&\code{["i", c, a, h]}, with \(c\not\equiv0\)\\
\(s=t\)&\code{["=", s, t]}\\
\(\Q_yB\)&\code{["q", B]}\\
\(\D_{x,y}B\)&\code{["d", B]}\\\bottomrule
\end{tabular}
\end{center}
A certificate has exactly four keys:
\code{variables}, \code{context}, \code{goal}, \code{proof}.
It specifies the number of free numerical variables, assumptions,
claimed formula, and finite proof tree. No executable code is accepted
inside a certificate.

The six proof tags are \code{hyp}, \code{succ}, \code{tr},
\code{qe}, \code{de}, \code{di}. The \code{tr} tag is X.
The checker infers each node's conclusion and compares the final
result with an independently validated goal. Every introduced term
and matrix is checked at the correct scope.

\subsection{The charged normalization}
The inherited traversal already keeps the iterator counter and
initializer at their surrounding scope and binds one variable only
in the body. After rebuilding a node, it now performs
\begin{lstlisting}[language=Python]
return a[2] if a[:2] == ['i', ['z']] else a
\end{lstlisting}
The public term validator additionally requires
\code{c != ['z']} at each iterator. The \code{zero} proof branch
is removed.

External syntax is validated before use. An out-of-scope variable
hidden in a supplied \(I(0,a;u.h)\) cannot be erased to evade validation:
the noncanonical node is rejected. The legacy proof compiler is
untrusted and validates its source certificate before normalization.
Native canonical certificates require no compiler.

\subsection{Soundness, completeness, and resources}
A valid finite six-rule proof tree passes the corresponding structural
checks. Conversely, each successful inference tag exhibits an instance
of its displayed rule, while scope and sort invariants are preserved
by hereditary substitution. Induction on the accepted tree therefore
gives a derivation in \(\Zsix\).

As in the baseline, this is the mathematical unbounded-resource
contract. A concrete Python process may reject sufficiently large
valid certificates through recursion or memory limits. These limits
do not create an acceptance shortcut.
Explicit \code{require} checks, rather than Python \code{assert},
keep validation active under \code{python3 -O}.

\section{Size comparison and limits}\label{sec:size}
The protocol is CPython 3.12.14: parse the full standalone source,
apply \code{ast.unparse}, append a newline, count UTF-8 bytes and
significant tokens, and count nodes from \code{ast.walk}.
Gzip uses level 9, fixed zero timestamp, and no filename, and is
informational.
\begin{center}
\small
\setlength{\tabcolsep}{5pt}
\begin{tabular}{lrrrrr}
\toprule
Complete tagged checker & Rules & Tokens & AST & Bytes & Gzip-9\\
\midrule
Raw, two introductions & 8 & 1,084 & 1,141 & 3,322 & 984 \\
Normal, two introductions & 7 & 1,073 & 1,136 & 3,272 & 996 \\
Raw, single introduction & 7 & 1,024 & 1,080 & 3,145 & 964 \\
Normal, single introduction & 6 & 1,013 & 1,075 & 3,095 & 977 \\
\bottomrule
\end{tabular}
\end{center}


The conservative comparison is the equally factored raw seven-rule
control. The final six-rule checker saves 50 bytes (about \(1.59\%\)),
11 tokens, and five AST nodes. The two-introduction normal checker
also saves 50 bytes against its matched eight-rule raw control.
Gzip size increases; it is reported rather than minimized.

All normal-form validation, smart reconstruction, logical rules,
scope handling, JSON parsing, final goal comparison, and process
accept/reject code are charged. Translators and test runners are
untrusted and unnecessary for verifying native canonical certificates.

\subsection{A rejected apparent saving}
We also examined common-source congruence:
\[
                      \frac{s=t\qquad s=u}{h(t)=h(u)}.
\]
It is equivalent to ordinary equality transport. An apparent 36-byte
saving disappeared after identical factoring of the existing
one-context control: both measured 3,322 bytes, 1,084 tokens, and
1,141 AST nodes. That refactoring is not counted as a mathematical
improvement. The final system retains the existing one-context rule.

\subsection{The separate compact frontier}
The 521-byte arity-coded checker has a different representation,
goal contract, and successful-trace argument, with guarded public
formulas and auxiliary resources. The tagged checker explicitly
validates all term and formula constructors. No reduction of the
521-byte architecture is established here. An absolute minimum
would additionally require a fixed description language and an
optimality proof; this report supplies neither.

\section{Verification and formalization scope}\label{sec:verification}
The recorded run has \textbf{500 process checks}: 400 expected
acceptances and 100 expected rejections, with \textbf{zero mismatches},
including runs with ordinary Python and with \code{-O}.
All 64 inherited positive certificates translate successfully to both
normal calculi.

The syntactic suite records 4,832 substitution-commutation checks;
462 checks each of idempotence, canonical scope preservation and
shift commutation; and 400 checks each of zero erasure, successor
computation preservation and formula-root preservation.
The exact inputs, tests and results are included in
\code{code/verify.py} and \code{evidence/verification.json}.


The suite includes genuine equation induction, QE and DE, compound
introduction, transport under binders, and the missing atomic case
introduction at literal-zero witness. That last case exercises
\(K=I(S0,0;u.0)\). A further test uses the replaced local assumption
under a nested DE eigenvariable, checking capture-avoiding weakening.

Negative cases address mismatched equality instances, forbidden
compound progress, escaping variables, malformed tags, scope errors,
and noncanonical iterators in goals, contexts, bodies, templates,
seeds, and counters.

The new 255-line \code{formal/ZeroNormalization.lean} compiles with
Lean 4.32.0 and Mathlib revision
\code{81a5d257c8e410db227a6665ed08f64fea08e997}. It proves:
\begin{itemize}
\item normal-form characterization, idempotence, renaming and lifting;
\item the full term and formula substitution-coherence laws;
\item hereditary substitution composition;
\item \(T_{\rm eq}\vdash t=\nf(t)\), using iterator-body congruence;
\item the ordinary formula equivalence and both raw native formula
return derivations, including \(\Q\) and \(\D\).
\end{itemize}
Its substitution and ordinary correctness theorems depend only on
Lean's \code{propext} and \code{Quot.sound}.
The native formula return theorem additionally inherits
\code{Classical.choice} from the existing native completeness proof.
There is no \code{sorryAx} or added axiom. These are metatheoretic
proof-assistant dependencies; the object logic remains intuitionistic.

The pinned native/iPRA dependency chain was rebuilt. The existing
one-context modules were also checked against that baseline and are
included as attributed references. One inherited namespace reference,
\code{WPrf.eqSymm}, was corrected to \code{WKernel.eqSymm}; the exact
compatibility patch is recorded separately and the frozen source is retained.

\textbf{Formalization boundary.}
The new six-rule derivability relation, its entire proof-tree
equivalence, and the Python checker are not asserted to be formally
verified in Lean. They have the explicit written proof above and
executable proof-translation tests. Lean verifies the new normalization
and return bridge; the inherited compiled sources verify the iPRA
bridge. Tests alone are not used as a universal equivalence proof.

\section{Reproduction and contents}
\begin{center}
\begin{tabular}{p{0.32\linewidth}p{0.58\linewidth}}
\toprule
Location&Purpose\\\midrule
\code{paper/report.tex}, \code{paper/report.pdf}
 & Complete report and compiled version.\\
\code{code/zero_six.py}
 & Main standalone checker.\\
\code{code/zero_seven.py}
 & Two-introduction bridge checker.\\
\code{code/}
 & Matched controls, translators, fixtures, examples, and verification.\\
\code{formal/ZeroNormalization.lean}
 & Compiled normalization and internal return bridge.\\
\code{reference/ConciseIPRA/}
 & Frozen Lean dependencies and one-context modules.\\
\code{reference/prior/}
 & Explicit antecedents for definitional zero computation.\\
\code{evidence/}
 & Metrics, verification, compilation records, provenance and checksums.\\
\bottomrule
\end{tabular}
\end{center}
The README gives exact commands. Python checks use only the standard
library. LaTeX uses ordinary packages. Lean reproduction requires
the stated compiler and pinned Mathlib.

A further step would formalize the normalized proof relation and
both proof translations, then examine whether the reduction survives
a fully charged adaptation of the much smaller arity codec.
The delivered result is the explicit six-rule theory and measured
direct-checker reduction.

\appendix
\section{The inherited iPRA interface}
For each fixed primitive recursive definition, iPRA has a function
symbol. Initial functions are zero, successor, and projections.
Composition introduces
\(f(\vec x)=g(h_1(\vec x),\ldots,h_k(\vec x))\).
Primitive recursion introduces
\[
 f(\vec x,0)=g(\vec x),\qquad
 f(\vec x,Sn)=h(\vec x,n,f(\vec x,n)).
\]
Successor injectivity, \(Sx=0\to\bot\), and quantifier-free induction
are added to intuitionistic first-order logic with equality.

The native-to-iPRA interpretation compiles iteration to primitive
recursion whose step ignores the counter. The reverse scalar
interpretation uses pairing to carry counter and state and thereby
simulates arbitrary primitive recursion with unary-state iteration.
The inherited bridge proves both composites for arbitrary formulas.
This report changes neither that compiler nor its numerical domain.

\section{A scope example}
The raw term
\[
                       I(n,a;u.I(0,u;v.0))
\]
normalizes to \(I(n,a;u.u)\) if the outer counter is not literally zero.
In the tagged syntax the inner initializer \(u\) was never under the
inner binder \(v\); it retains the correct index after erasure.
This avoids the dummy-slot lowering cost in the older arity codec.

By contrast \(I(S0,0;u.0)\) does not normalize to zero. Its equality
to zero is a theorem of IS. Keeping this distinction enables the
single-introduction argument.

\begin{thebibliography}{99}\footnotesize
\setlength{\itemsep}{1pt}\setlength{\parskip}{0pt}\setlength{\parsep}{0pt}
\bibitem{readme}
Vladimir Reshetnikov, \emph{ProveIt: Concise presentations of intuitionistic
primitive recursive arithmetic}, commit
\code{ecb3b847290c33b1ef7f04571cac63e97c73a54e}.
\url{https://github.com/VladimirReshetnikov/ProveIt/tree/ecb3b847290c33b1ef7f04571cac63e97c73a54e/Logic/PrimitiveRecursiveArithmetic/ConciseIPRA}.
\bibitem{native}
Same repository and commit, \code{Lean/ConciseIPRA/Native.lean},
\code{NativeEq.lean}, \code{NativeGate.lean}, \code{NativeLogic.lean}.
Native rules, equality, selectors and logical macros.
\bibitem{oneintro}
Same, \code{Lean/ConciseIPRA/NativeOneIntro.lean}.
The witness-selected introduction and \code{wprf_iff_nprf}.
\bibitem{onecontext}
Same, \code{Research/one_context/}.
The crossed one-context rule, tagged checkers, and Lean comparison modules.
\bibitem{priorzero}
Same, \code{Research/prevalidated_elimination/theory/research.tex},
subsection \emph{Making zero computation definitional}.
\bibitem{uniformzero}
Same, \code{Research/uniform_zero/report.tex},
question on definitional zero computation with a different binder layout.
\bibitem{roundtrip}
Same, \code{Lean/ConciseIPRA/RoundTripScalar.lean},
\code{Q.iter_body_congr} and \code{Q.iter_congr}.
\bibitem{nativebridge}
Same, \code{Lean/ConciseIPRA/NativeBiInterpretation.lean} and dependency
chain, including \code{BiInterpretation.lean}.
Both interpretations, formula round trips and equiconsistency.
\bibitem{frontier}
Same, \code{Research/kernel_refinement/} and \code{Research/compact_update/}.
The 521-byte arity-coded frontier and its stated contract.
\bibitem{this}
This package, \code{code/}, \code{formal/}, \code{evidence/}.
Reproducible metrics, translations, tests, and compiled normalization.
\end{thebibliography}
\end{document}
